%! Characteristics of Polynomial Functions — Unit 3, Lesson 3.0 — Slides (Beamer)
\documentclass[aspectratio=169, 11pt]{beamer}
\usepackage{algebra2-beamer}   % provides \forestheader{} and \sectionlabel[color]{}
% Standard coordinate grid + axes: {xmin}{xmax}{ymin}{ymax}
\newcommand{\lgrid}[4]{%
  \draw[step=1, linegray!45, thin] (#1,#3) grid (#2,#4);
  \draw[->, thick, charcoal] (#1,0)--(#2,0) node[below right, font=\tiny]{$x$};
  \draw[->, thick, charcoal] (0,#3)--(0,#4) node[left, font=\tiny]{$y$};}

\begin{document}

% TITLE SLIDE (CourseName is not defined in beamer, so the name is literal)
{
\setbeamercolor{background canvas}{bg=forest}
\begin{frame}[plain]
  \vspace{2.8cm}\hspace{0.55cm}%
  \begin{minipage}{0.9\textwidth}
    {\color{goldacc}\bfseries\small LESSON 3.0}\\[0.5cm]
    {\color{white}\bfseries\LARGE Characteristics of Polynomial Functions}\\[1.2cm]
    {\color{forestmid}\small Algebra 2: Shepherd~$\cdot$~Unit 3: Polynomial Functions}
  \end{minipage}
\end{frame}
}

% HOOK
\begin{frame}[t, plain]
  \forestheader{A parabola turns once --- this turns twice}
  \sectionlabel{Hook}
  \begin{columns}[T]
    \begin{column}{0.54\textwidth}
      \[ g(x) = x^3 - 3x + 2 = (x-1)^2(x+2) \]
      \begin{itemize}\setlength{\itemsep}{4pt}
        \item Top of the hill / bottom of the valley? $\to$ \textbf{relative extrema}
        \item Far-left / far-right arms? $\to$ \textbf{end behavior}
        \item Bounces at $x=1$, cuts through at $x=-2$? $\to$ \textbf{multiplicity}
      \end{itemize}
    \end{column}
    \begin{column}{0.44\textwidth}
      \centering
      \begin{tikzpicture}[scale=0.42]
        \lgrid{-3}{3}{-3}{5}
        \begin{scope}
          \clip (-3,-3) rectangle (3,5);
          \draw[very thick, forest, domain=-2.15:2.1, samples=90, smooth]
            plot (\x,{(\x)*(\x)*(\x) - 3*(\x) + 2});
        \end{scope}
        \fill[charcoal] (-1,4) circle (4pt); \fill[charcoal] (1,0) circle (4pt);
        \fill[charcoal] (-2,0) circle (4pt); \fill[charcoal] (0,2) circle (4pt);
      \end{tikzpicture}
    \end{column}
  \end{columns}
  \vspace{2pt}
  \begin{block}{Today}
    Read \emph{every} characteristic of a polynomial straight off its graph.
  \end{block}
\end{frame}

% DEFINITION: degree & leading coefficient
\begin{frame}[t, plain]
  \forestheader{Polynomials: degree and leading coefficient}
  \sectionlabel{Definition}
  \begin{columns}[T]
    \begin{column}{0.55\textwidth}
      A \textbf{polynomial} is a sum of terms $a_kx^k$ with \emph{whole-number} exponents ---
      a \emph{smooth}, \emph{continuous} curve.
      \[ f(x)=a_nx^n+\dots+a_1x+a_0 \]
      \begin{itemize}\setlength{\itemsep}{4pt}
        \item \textbf{Degree:} the highest power $n$.
        \item \textbf{Leading coefficient:} $a_n$, in standard form.
      \end{itemize}
    \end{column}
    \begin{column}{0.42\textwidth}
      \begin{block}{Not polynomials}
        $\dfrac{1}{x}=x^{-1}$, \quad $\sqrt{x}=x^{1/2}$ \\[3pt]
        {\footnotesize exponents must be whole numbers.}
      \end{block}
      \vspace{2pt}
      {\footnotesize For $g(x)=x^3-3x+2$: degree $3$, leading coefficient $1$.}
    \end{column}
  \end{columns}
\end{frame}

% END BEHAVIOR: the two-question rule
\begin{frame}[t, plain]
  \forestheader{End behavior: only two questions}
  \sectionlabel[goldacc]{Even/odd degree \& sign of $a$}
  \begin{center}
  \begin{tabular}{cccc}
  \begin{tikzpicture}[scale=0.34]
    \draw[->, charcoal] (-2,0)--(2,0); \draw[->, charcoal] (0,-0.6)--(0,4);
    \draw[very thick, forest, domain=-1.7:1.7, samples=40, smooth] plot (\x,{(\x)*(\x)});
  \end{tikzpicture} &
  \begin{tikzpicture}[scale=0.34]
    \draw[->, charcoal] (-2,0)--(2,0); \draw[->, charcoal] (0,-4)--(0,0.6);
    \draw[very thick, forest, domain=-1.7:1.7, samples=40, smooth] plot (\x,{-1*(\x)*(\x)});
  \end{tikzpicture} &
  \begin{tikzpicture}[scale=0.34]
    \draw[->, charcoal] (-2,0)--(2,0); \draw[->, charcoal] (0,-3.2)--(0,3.2);
    \draw[very thick, forest, domain=-1.45:1.45, samples=50, smooth] plot (\x,{(\x)*(\x)*(\x)});
  \end{tikzpicture} &
  \begin{tikzpicture}[scale=0.34]
    \draw[->, charcoal] (-2,0)--(2,0); \draw[->, charcoal] (0,-3.2)--(0,3.2);
    \draw[very thick, forest, domain=-1.45:1.45, samples=50, smooth] plot (\x,{-1*(\x)*(\x)*(\x)});
  \end{tikzpicture} \\[2pt]
  {\footnotesize even, $a>0$} & {\footnotesize even, $a<0$} & {\footnotesize odd, $a>0$} & {\footnotesize odd, $a<0$} \\
  {\footnotesize both up} & {\footnotesize both down} & {\footnotesize L$\downarrow$ R$\uparrow$} & {\footnotesize L$\uparrow$ R$\downarrow$} \\
  \end{tabular}
  \end{center}
  \begin{itemize}\setlength{\itemsep}{3pt}
    \item \textbf{Even} degree $\Rightarrow$ arms the \emph{same} way; \textbf{odd} degree $\Rightarrow$ arms \emph{opposite}.
    \item Sign of $a$ sets the \emph{right} arm: $+\Rightarrow$ up, $-\Rightarrow$ down.
  \end{itemize}
\end{frame}

% TURNING POINTS & EXTREMA
\begin{frame}[t, plain]
  \forestheader{Turning points: relative vs.\ absolute}
  \sectionlabel[goldacc]{$g(x)=x^3-3x+2$}
  \begin{columns}[T]
    \begin{column}{0.44\textwidth}
      \centering
      \begin{tikzpicture}[scale=0.42]
        \lgrid{-3}{3}{-3}{5}
        \begin{scope}
          \clip (-3,-3) rectangle (3,5);
          \draw[very thick, forest, domain=-2.15:2.1, samples=90, smooth]
            plot (\x,{(\x)*(\x)*(\x) - 3*(\x) + 2});
        \end{scope}
        \fill[charcoal] (-1,4) circle (4pt) node[above right, font=\tiny]{$(-1,4)$};
        \fill[charcoal] (1,0) circle (4pt) node[below right, font=\tiny]{$(1,0)$};
      \end{tikzpicture}
    \end{column}
    \begin{column}{0.56\textwidth}
      \footnotesize
      \begin{itemize}\setlength{\itemsep}{4pt}
        \item A \textbf{turning point} is where the curve changes rise$\leftrightarrow$fall; degree $n$ has \emph{at most} $n-1$.
        \item \textbf{Relative (local) max} $(-1,4)$, \textbf{relative min} $(1,0)$ --- highest/lowest \emph{nearby}.
        \item \textbf{Absolute (global)} max/min = highest/lowest of the \emph{whole} function.
        \item Odd degree $\Rightarrow$ arms split to $\pm\infty$ $\Rightarrow$ \emph{no} absolute max or min.
      \end{itemize}
    \end{column}
  \end{columns}
\end{frame}

% ZEROS & MULTIPLICITY
\begin{frame}[t, plain]
  \forestheader{Multiplicity: cross vs.\ touch}
  \sectionlabel{Zeros}
  \begin{columns}[T]
    \begin{column}{0.44\textwidth}
      \centering
      \begin{tikzpicture}[scale=0.42]
        \lgrid{-3}{3}{-3}{5}
        \begin{scope}
          \clip (-3,-3) rectangle (3,5);
          \draw[very thick, forest, domain=-2.15:2.1, samples=90, smooth]
            plot (\x,{(\x)*(\x)*(\x) - 3*(\x) + 2});
        \end{scope}
        \fill[charcoal] (-2,0) circle (4pt) node[above left, font=\tiny]{cross};
        \fill[charcoal] (1,0) circle (4pt) node[above right, font=\tiny]{touch};
      \end{tikzpicture}
    \end{column}
    \begin{column}{0.56\textwidth}
      \[ g(x)=(x-1)^2(x+2) \]
      \begin{itemize}\setlength{\itemsep}{5pt}
        \item \textbf{Multiplicity} $=$ the exponent of the factor.
        \item $x=-2$: multiplicity $1$ (\textbf{odd}) $\Rightarrow$ graph \textbf{crosses}.
        \item $x=1$: multiplicity $2$ (\textbf{even}) $\Rightarrow$ graph \textbf{touches} (bounces).
      \end{itemize}
      \begin{block}{Remember}
        \textbf{Odd crosses, even touches.}
      \end{block}
    \end{column}
  \end{columns}
\end{frame}

% SYMMETRY
\begin{frame}[t, plain]
  \forestheader{Symmetry: even vs.\ odd}
  \sectionlabel[goldacc]{The $f(-x)$ test}
  \begin{columns}[T]
    \begin{column}{0.5\textwidth}
      \centering
      \begin{tikzpicture}[scale=0.38]
        \lgrid{-2}{2}{-1}{5}
        \begin{scope}
          \clip (-2,-1) rectangle (2,5);
          \draw[very thick, forest, domain=-1.5:1.5, samples=60, smooth] plot (\x,{(\x)^4});
        \end{scope}
        \fill[charcoal] (-1,1) circle (3pt); \fill[charcoal] (1,1) circle (3pt);
      \end{tikzpicture}\\
      {\footnotesize $y=x^4$: \textbf{even}, $f(-x)=f(x)$, $y$-axis mirror}
    \end{column}
    \begin{column}{0.5\textwidth}
      \centering
      \begin{tikzpicture}[scale=0.38]
        \lgrid{-2}{2}{-4}{4}
        \begin{scope}
          \clip (-2,-4) rectangle (2,4);
          \draw[very thick, forest, domain=-1.55:1.55, samples=60, smooth] plot (\x,{(\x)*(\x)*(\x)});
        \end{scope}
        \fill[charcoal] (1,1) circle (3pt); \fill[charcoal] (-1,-1) circle (3pt);
      \end{tikzpicture}\\
      {\footnotesize $y=x^3$: \textbf{odd}, $f(-x)=-f(x)$, $180^\circ$ about origin}
    \end{column}
  \end{columns}
  \vspace{3pt}
  \begin{center}\footnotesize Only even powers $\Rightarrow$ even; only odd powers $\Rightarrow$ odd; a mix $\Rightarrow$ neither.\end{center}
\end{frame}

% ROADMAP
\begin{frame}[t, plain]
  \forestheader{Where Unit 3 goes next}
  \sectionlabel{Connections}
  \textbf{Today:} read degree/leading coefficient, end behavior, turning points, relative/absolute
  extrema, zeros \& multiplicity, and symmetry from a polynomial graph.\\[8pt]
  \begin{itemize}\setlength{\itemsep}{3pt}
    \item \textbf{3.1} Operations with polynomials (add, subtract, multiply).
    \item \textbf{3.2} Advanced factoring (grouping, sum \& difference of cubes).
    \item \textbf{3.3} Dividing polynomials; Remainder \& Factor Theorems.
    \item \textbf{3.4--3.5} Zeros: Rational Root Theorem; complex zeros \& the FTA.
    \item \textbf{3.6} Graphing polynomial functions \& modeling.
  \end{itemize}
\end{frame}

\end{document}
